%% kapa: exam=1 round=27 year=2016 period=1 form=A subject=real_estate subject_ko=부동산학원론 seq=37 number=117
%% source: Q-Net 감정평가사 기출문제 (cst00309) artlSeq=5169402 file=extracted/1008/5169402/제27회 감정평가사 1차 문제/제27회 감정평가사 1교시문제(A형).pdf
%% conversion: kordoc 4.24.51 -> tools/gen1.py (auto); review=proofread-sonnet; flags=none
\begin{question}{117}
부동산마케팅에 관한 설명으로 옳은 것은?
\begin{choices}
  \item 표적시장(Target market)은 목표시장에서 고객의 욕구를 파악하여 경쟁 제품과 차별성을 가지도록 제품 개념을 정하고 소비자의 지각 속에 적절히 위치시키는 것이다.
  \item 포지셔닝(Positioning)은 세분화된 시장 중 가장 좋은 시장기회를 제공해 줄 수 있는 특화된 시장이다.
  \item 4P에 의한 마케팅 믹스 전략의 구성요소는 제품(Product), 유통경로(Place), 판매 촉진(Promotion), 포지셔닝(Positioning)이다.
  \item STP란 시장세분화(Segmentation), 표적화(Targeting), 가격(Price)을 표상하는 약자이다.
  \item 고객점유 마케팅 전략은 AIDA(Attention, Interest, Desire, Action) 원리를 적용하여 소비자의 욕구를 충족시키기 위한 마케팅 전략이다.
\end{choices}
\end{question}
\answer{5}
